\chapter{Inductive Bias, Invariance, and Equivariance}
\label{ch:inductive-bias}
\minitoc

The preceding analysis in Section~\ref{sec:mlp-fails} introduced two concepts
in passing — the \textit{inductive bias} of a model and the requirement that a
graph model be \textit{invariant} to node ordering — without developing either.
This chapter develops both. It first makes precise what an inductive bias is
and why successful generalization requires one, reading the inductive biases of
successful architectures as evidence that each encodes the structure of its own
domain. It then identifies the inductive bias appropriate to graphs and gives
the two symmetry conditions that express it formally.

\section{Inductive Bias in Neural Architectures}
\label{sec:inductive-bias}

A model trained on a finite set of examples must, to be useful, make
predictions on examples it has never seen. The training data alone cannot
determine those predictions: infinitely many functions are consistent with any
finite sample. A model generalizes only because it brings assumptions of its
own about which functions are plausible before any data is observed.
The set of such assumptions is the \textit{inductive bias} of the
model~\cite{goodfellow2016}. Generalization is therefore not possible without
an inductive bias, and succeeds when that bias matches the structure of the
problem.

Deep learning architectures differ primarily in the inductive biases their
structure encodes. Three of them recur throughout 
this work: the multilayer perceptron, the convolutional neural network (CNN), 
and the recurrent neural network (RNN). From this point onward, these architectures 
are referred to exclusively by their acronyms, both in this chapter and 
throughout the remainder of the document.

The MLP, whose layers connect every input to every output, assumes that any
input coordinate may interact with any other, that no particular structure 
among the coordinates is presumed, and that examples are treated 
independently. This is the appropriate bias when the input is an
unstructured vector — such as tabular data — but the wrong bias for 
data with an underlying structure.

The CNN and the RNN illustrate how a well-matched bias is built into
structure. The CNN, designed for data arranged on a regular grid such as the
pixels of an image, applies the same local filter at every position. Two
assumptions are encoded in this design. The first is \textit{locality}: a pixel
is informative about its immediate neighbors, so the filter reads only a small
spatial window rather than the whole image. The second is that a pattern is
meaningful regardless of where in the image it appears, so the same filter is
reused at every position; a feature shifted across the grid produces the same
feature shifted in the output. These assumptions make the CNN well suited to 
natural images. The RNN, designed for sequences,
encodes a different bias: it processes elements in order, applying the same
transition at every step. This assumes that the order of the
elements is meaningful and that the same transformation can be reused across
all positions in the sequence. These assumptions are appropriate for text and time series.

The same principle appears in all three. Each architecture succeeds when its
inductive bias matches the structure of the data, and the match is built into how
the architecture is wired rather than learned from data. The question
this chapter must answer follows immediately: what is the structure of a graph,
and what inductive bias matches that structure.

\section{The Relational Inductive Bias}
\label{sec:relational-bias}

The defining structure of a graph is relational. An entity is not described in
isolation but through its connections to other entities: in a citation network
the subject of a paper is reflected in the papers it cites, and in a molecule
the role of an atom is determined by the atoms bonded to it. The inductive bias
appropriate to graphs is therefore the assumption that a node is characterized
by its neighborhood --- that the information relevant to a node resides in the node itself,
in the nodes adjacent to it, and in the structure connecting them. This assumption is
the \textit{relational inductive bias}~\cite[-7.5pc]{battaglia2018}.

The relational inductive bias is the graph counterpart of the locality that the CNN encodes,
and the comparison is exact in one respect and instructive in its difference. Both assume
that an entity is informed by what lies near it rather than by everything at once,
and both read from a local region rather than the whole input. The difference lies in what
counts as local. On a grid, locality is geometric: a pixel's neighbors are the pixels at fixed
offsets around it, the same offsets everywhere, so a local window is defined once and
applied uniformly. On a graph there are no fixed offsets and no uniform window.
A node's neighborhood is whatever the edges make it, varying in size from one node to
the next. The relational inductive bias retains locality while discarding the regular geometry
that made it simple to express on a grid; locality on a graph is
defined by adjacency, not by position.

Adjacency, however, fixes only which nodes are neighbors, not the order in which
they are listed. A neighborhood is a set, and a set has no first element. The
relational inductive bias says where a node's information comes from; it does
not, and must not, depend on the order in which that information is enumerated,
because the graph supplies no such order. Making this requirement precise calls
for the concepts of the next two sections.

\section{Invariance and Equivariance}
\label{sec:invariance-equivariance}

The requirement that a model not depend on an arbitrary choice in its input is
made precise by the notions of \textit{invariance} and \textit{equivariance}. 
Both describe how a function responds to transformations of its input, 
and they are defined for any function and any transformation 
before being specialized to graphs.

Let $f$ be a function and let $g$ denote a transformation applied to its input.
The function $f$ is \textit{invariant} to $g$ when transforming the input leaves
the output unchanged:
\begin{equation}
  f(g \cdot \mathbf{x}) = f(\mathbf{x}).
  \label{eq:invariance}
\end{equation}
The transformation is applied to the input, and the output does not move. The
function $f$ is \textit{equivariant} to $g$ when transforming the input
transforms the output in the same way:
\begin{equation}
  f(g \cdot \mathbf{x}) = g \cdot f(\mathbf{x}).
  \label{eq:equivariance}
\end{equation}
The transformation is applied to the input, and
the output moves in step with the input.

The distinction is between discarding a transformation and preserving it.
Invariance identifies transformed and untransformed inputs by mapping 
them to the same output, so the transformation is discarded by $f$. 
Equivariance preserves the effect of the transformation:
it survives the application of $f$ and acts on the result, so applying $f$
and then transforming gives the same answer as transforming and then applying
$f$. Which of the two is required is not a matter of preference but of the task:
when the quantity to be predicted is itself unaffected by the transformation,
$f$ must be invariant; when the quantity to be predicted is attached to parts of
the input that the transformation moves, $f$ must be equivariant, so that the
prediction moves with the part it describes.

\section{Permutation Invariance and Equivariance on Graphs}
\label{sec:permutation}

The transformation that matters for graphs is the reordering of nodes. As
established in Chapter~\ref{ch:graph-theory}, the nodes of a graph carry no
intrinsic order; an indexing must be chosen to write the adjacency matrix
$\mathbf{A}$ and the feature matrix $\mathbf{X}$, and any of the $n!$ possible
indexings describes the same graph. A reordering of the nodes is represented by
a \textit{permutation matrix} $\mathbf{P} \in \{0,1\}^{n \times n}$, a square
matrix with exactly one entry equal to $1$ in each row and each column and zeros
elsewhere. An entry $P_{ij}=1$ indicates that the node originally at index $j$ is mapped to
the new index $i$. Multiplying by $\mathbf{P}$ rearranges the rows of a matrix 
according to this relabeling. A permutation matrix is orthogonal,
\begin{equation}
  \mathbf{P}\mathbf{P}^\top = \mathbf{I},
  \label{eq:permutation-orthogonal}
\end{equation}
which states that the reordering is invertible and that its inverse is its
transpose.

Applying the reordering $\mathbf{P}$ to a graph acts on both of its matrices.
The feature matrix becomes $\mathbf{P}\mathbf{X}$, which moves each node's
feature vector to the row assigned by the reordering. The adjacency matrix
becomes $\mathbf{P}\mathbf{A}\mathbf{P}^\top$, with $\mathbf{P}$ reordering its
rows and $\mathbf{P}^\top$ its columns, so that the entry recording an edge
between two nodes follows them to their new positions. The pair
$(\mathbf{P}\mathbf{A}\mathbf{P}^\top,\, \mathbf{P}\mathbf{X})$ is the same graph
as $(\mathbf{A}, \mathbf{X})$ under a different indexing; no edge has been added
or removed, and no feature has changed.

A function defined on graphs must respect this fact, and which condition it must
satisfy depends on the level of the task. For a task whose output describes the
graph as a whole --- a single prediction for the entire graph, as in graph
classification --- the output must not depend on the indexing at all. Such a
function $f$ must be \textit{permutation invariant}:
\begin{equation}
  f(\mathbf{P}\mathbf{A}\mathbf{P}^\top,\, \mathbf{P}\mathbf{X})
    = f(\mathbf{A}, \mathbf{X}).
  \label{eq:permutation-invariance}
\end{equation}
Reordering the nodes leaves the graph-level prediction unchanged.
For a task whose output is attached to the individual 
nodes --- one prediction per node, as in node classification --- the output 
must follow the nodes when they are reordered.
In this case, $f$ must be \textit{permutation equivariant}:
\begin{equation}
  f(\mathbf{P}\mathbf{A}\mathbf{P}^\top,\, \mathbf{P}\mathbf{X})
    = \mathbf{P}\, f(\mathbf{A}, \mathbf{X}).
  \label{eq:permutation-equivariance}
\end{equation}
Reordering the nodes reorders the per-node outputs in exactly the same way: the
prediction for a given node is the same prediction, carried to that node's new
position~\cite[5.7pc]{hamilton2020}. Figure~\ref{fig:permutation-equivariance}
illustrates the two conditions on a single graph drawn under two orderings.

\begin{figure*}[t]
  \centering
  \includegraphics{permutation-equivariance}
  \caption[Permutation invariance and equivariance on a graph.]{%
    A single graph shown under two node orderings related by a permutation
    $\mathbf{P}$. Left: the original indexing, with a per-node output attached
    to each node and a single graph-level output computed from the whole graph.
    Right: the reordered indexing, with the graph-level output unchanged
    (invariance), while the per-node outputs appear in the permuted positions
    (equivariance).
  }
  \label{fig:permutation-equivariance}
\end{figure*}

These two conditions, together with the relational inductive bias, 
determine what a layer operating on graphs
must be. The relational inductive bias fixes where each node's information comes
from. Permutation equivariance fixes how the layer must treat
the indexing. This is the design the next chapter develops: the message-passing 
framework gives these requirements a concrete and general form.